\begin{enumerate}
    \item Soit $G$ un graphe connexe. Il existe au moins deux composantes connexes dans $G$. Retirer une arête à $G$ revient à retirer une arête à l'une de ces composantes connexes, au risque~:
    \begin{itemize}
        \item ou bien de la diviser en deux nouvelles composantes connexes~;
        \item ou bien de garder autant de composantes connexes.
    \end{itemize}
    Ainsi, il existe au moins autant de composantes connexes à $\Psi(G)$, c'est-à-dire au moins~$2$. D'où $\Psi(G)$ n'est pas connexe.
    \item Si $G$ est acyclique, alors on peut affirmer qu'il manque une ou des arêtes à $G$ pour le rendre cyclique. De fait, en retirer une ne va pas arranger les choses. $\Psi(G)$ est donc acyclique.
    \item Le graphe $G$ ci-dessous est cyclique. Pourtant $\Psi(G)$ est acyclique~:
    \begin{tikzpicture}
        \node[draw,circle] (A) at(0,0) {$A$};
        \node[draw,circle] (B) at(0,0) {$B$};
        \node[draw,circle] (B) at(0,0) {A};
    \end{tikzpicture}
\end{enumerate}
